% [AUTHOR REWRITE] empirical reading of the proposition
The proposition gives the two segments one reading. In apartments the non-disclosure premium, \attDthreeApt{} log points, is about twice the energy-efficiency premium, \attBthreeApt{}: the effective belief is \attRatioApt{}, and silent listings sell above verified efficient ones, which sell above inefficient ones. In houses the effective belief is \attRatioHouse{} and silence sells between the two ratings. These orderings are read directly from the two estimated premia. Which case of Corollary~\ref{cor:att_ordering} each segment falls in also depends on the normalisation of the default ($C^{\rm def}=0$): at the estimated attention apartments are in the first case ($\hat\psi_\delta=\attPsiApt{}$ against a threshold of \attPsiBarOneApt{}), so silence outprices a verified efficient rating whatever the hiding cost, and houses are in the second ($\attPsiHouse{}$ against \attPsiBarThetaHouse{}). The model with one $\psi$ per segment therefore predicts the first ordering in every apartment market and the second in every house market. The market-level estimates are noisy and $\psi$ is common within a segment, so the counts differ: silence sells above a verified efficient rating in \attNDgtBApt{} of \attNMarketsApt{} apartment markets and \attNDgtBHouse{} of \attNMarketsHouse{} house markets, between the two displayed ratings in \attNBetweenApt{} and \attNBetweenHouse{}, and below a displayed inefficient rating in \attRawNNegApt{} and \attRawNNegHouse{}. The last cannot arise in Partial Pooling and reflects sampling error or heterogeneity that we do not model. The payoff ranking among informed sellers is stated at the hiding cost that reproduces the observed hiding rate, which at the aggregate and in Partial Pooling is the premium itself, $c_0=\hat\Delta$ (Proposition~\ref{prop:att_identification}(iv)); the price ordering needs only $\hat\Delta$ against $\hat\beta_1$. Read through Partial Pooling, an uninformed apartment seller earns \attDthreeApt{} log points over a displayed inefficient rating, an informed high-quality seller \attBthreeApt{}, and an informed low-quality seller nothing, whether she hides or not. Silence paying more than a verified rating does not make high-quality sellers hide, which is an implication of Partial Pooling and not a separate finding: their gain, \attExcessApt{} log points, falls short of the hiding cost \attDthreeApt{} that keeps low-quality sellers indifferent. The uninformed seller's rent is a property of the maintained assumptions, that uninformed sellers have no strategy, bear no hiding cost and are exogenous in number, and not a finding.
